% Generated from locale/tr/content/first-order-logic/syntax-and-semantics/unique-readability.tex; do not hand-edit.


\olfileid[tr]{fol}{syn}{unq}

\olsection{Tek Türlü Okunabilirlik}

\begin{explain}
!!{formula}s için verdiğimiz tanım, her !!{formula} açısından \emph{tek türlü
okunabilirliği} güvence altına alır; başka bir deyişle, !!{formula}s için
verdiğimiz oluşum kurallarına göre onu kurmanın ve ``yorumlamanın'' özünde
yalnızca bir yolu vardır. Böyle olmasaydı belirsiz !!{formula}s ortaya
çıkardı; yani birden çok okuması ya da yorumu olan !!{formula}s söz konusu olurdu---ki
bunu istemediğimiz açıktır. Daha da önemlisi, bu özellik olmasaydı
vereceğimiz tanımların ve ispatların çoğu işlemezdi.

Bunu açıklığa kavuşturmanın belki de en iyi yolu, !!{formula}s için tek
türlü okunabilirliği güvence altına almayan kötü oluşum kuralları
vermiş olsaydık ne olacağını görmektir. Örneğin bağlaçlara ilişkin
oluşum kurallarında parantezleri unutmuş ve şuna izin vermiş
olabilirdik:
\begin{quote}
$!A$ ve $!B$ !!{formula}s ise $!A \lif !B$ de formüldür.
\end{quote}
Atomik bir $!D$ formülünden başlayarak bu kural $!D \lif !D$'yi
kurmamıza izin verirdi. Bundan ve $!D$'den de $!D \lif !D \lif !D$
elde ederdik. Fakat bunu yapmanın iki yolu vardır:
\begin{enumerate}
\item $!D$'yi $!A$, $!D \lif !D$'yi ise $!B$ olarak alırız.
\item $!A$'yı $!D \lif !D$, $!B$'yi ise~$!D$ olarak alırız.
\end{enumerate}
Buna karşılık $!D \lif !D \lif !D$ için ortaya çıkan !!{formula}
biçimini ``okumanın'' da iki yolu vardır. $!B$'nin $!D$, $!C$'nin $!D \lif !D$ olduğu $!B \lif !C$
biçimindedir; ama $!B$'nin $!D \lif !D$, $!C$'nin ise~$!D$ olduğu $!B
\lif !C$ biçiminde de \emph{bulunur}.

Böyle bir şey olursa tanımlarımız her zaman işlemez. Örneğin bir
formül için !!{main operator} kavramını tanımlarken şöyle deriz: $!B \lif !C$
biçimindeki bir formülde !!{main operator}, gösterilen $\lif$
oluşumudur. Fakat $!D \lif !D \lif !D$ formülünü yukarıda belirtilen iki
farklı yoldan $!B \lif !C$ ile eşleştirebiliyorsak, bir durumda
$\lif$'in ilk oluşumunu, öteki durumda ise ikinci oluşumunu !!{main
operator} olarak elde ederiz. Oysa !!{main operator} belirleme işleminin
!!{formula} üzerinde bir \emph{fonksiyon} olmasını amaçlarız; yani her !!{formula} için tam
olarak bir !!{main
  operator} oluşumu bulunmalıdır.
\end{explain}

\begin{lem}
Bir !!a{formula}~$!A$ içindeki sol ve sağ parantezlerin sayıları
eşittir.
\end{lem}

\begin{proof}
Bunu $!A$'nın kuruluşu üzerinde tümevarımla ispatlarız. Bunun için iki
şey gerekir: (a) Önce bütün atomik formüllerin söz konusu özelliğe sahip
olduğunu ispatlamalıyız (tümevarım tabanı). (b) Sonra, verilmiş
formüllerden yeni formüller kurduğumuzda, eski formüller bu özelliğe
sahipse yeni formüllerin de sahip olduğunu ispatlamalıyız.

$l(!A)$, $!A$'daki sol parantezlerin, $r(!A)$ ise sağ parantezlerin
sayısı olsun; benzer biçimde $l(t)$ ve $r(t)$ de $t$ terimindeki sol ve
sağ parantezlerin sayıları olsun.

\begin{prob}
    Her $t$ terimi için $l(t) = r(t)$ olduğunu ispatlayın.
\end{prob}

\begin{enumerate}
\tagitem{prvFalse}{\indcase{!A}{\lfalse}{$\indfrm$ $0$ sol ve $0$
    sağ parantez içerir.}}{}

\tagitem{prvTrue}{\indcase{!A}{\ltrue}{$\indfrm$ $0$ sol ve $0$
    sağ parantez içerir.}}{}

\item \indcase{!A}{\Atom{R}{t_1,\dots,t_n}}{$l(\indfrm) = 1 + l(t_1) +
  \dots + l(t_n) = 1 + r(t_1) + \dots + r(t_n) = r(\indfrm)$. Burada
  alıştırma olarak bırakılan, her $t$ terimi için $l(t) = r(t)$ olduğu
  olgusunu kullanırız.}

\item \indcase{!A}{\eq[t_1][t_2]}{$l(\indfrm) = l(t_1) + l(t_2) =
  r(t_1) + r(t_2) = r(\indfrm)$.}

\tagitem{prvNot}{\indcase{!A}{\lnot !B}{Tümevarım varsayımı uyarınca
    $l(!B) = r(!B)$'dir. Dolayısıyla $l(\indfrm) = l(!B) = r(!B) =
    r(\indfrm)$'dir.}}{}

\item \indcase{!A}{(!B \ast !C)}{Tümevarım varsayımı uyarınca $l(!B) =
  r(!B)$ ve $l(!C) = r(!C)$'dir. Dolayısıyla $l(\indfrm) = 1 + l(!B) + l(!C) =
  1 + r(!B) + r(!C) = r(\indfrm)$'dir.}

\tagitem{prvAll}{\indcase{!A}{\lforall[x][!B]}{Tümevarım
    varsayımı uyarınca $l(!B) = r(!B)$'dir. Dolayısıyla $l(\indfrm) = l(!B) = r(!B) =
    r(\indfrm)$'dir.}}{}

% Print case for \lexists if prvEx and not prvAll
\tagitem{prvAll,notprvEx}{}{\indcase{!A}{\lexists[x][!B]}{Tümevarım
    varsayımı uyarınca $l(!B) = r(!B)$'dir. Dolayısıyla $l(\indfrm) = l(!B) = r(!B) =
    r(\indfrm)$'dir.}}

% Just say similarly if prvEx and prvAll
\tagitem{notprvAll,notprvEx}{}{\indcase{!A}{\lexists[x][!B]}{Benzer biçimde.}}
\end{enumerate}
\end{proof}

\begin{defn}[Öz önek]
$!B$ simge dizisi, $!A$ simge dizisinin \emph{öz öneki}dir; bunun
anlamı, boş olmayan bir simge dizisinin $!B$'ye eklenmesiyle~$!A$ elde
edilmesidir.
\end{defn}

\begin{lem}\ollabel{lem:no-prefix}
$!A$ bir !!a{formula} ve $!B$, $!A$'nın öz önekiyse $!B$ bir
!!a{formula} değildir.
\end{lem}

\begin{proof}
Alıştırma.
\end{proof}

\begin{prob}
\olref[fol][syn][unq]{lem:no-prefix}'i ispatlayın.
\end{prob}

\begin{prop}
\ollabel{prop:unique-atomic}
$!A$ atomik bir !!{formula} ise aşağıdaki koşullardan yalnızca birini
sağlar.
\begin{enumerate}
\tagitem{prvFalse}{$!A \ident \lfalse$.}{}
\tagitem{prvTrue}{$!A \ident \ltrue$.}{}
\item $!A \ident \Atom{R}{t_1,\dots,t_n}$; burada $R$, $n$-yerli bir
  !!{predicate}; $t_1$, \dots, $t_n$ terimlerdir ve $R$,
  $t_1$, \dots, $t_n$ öğelerinin her biri tek türlü belirlenmiştir.
\item $!A \ident \eq[t_1][t_2]$; burada $t_1$ ile $t_2$ tek türlü
  belirlenmiş terimlerdir.
\end{enumerate}
\end{prop}

\begin{proof}
Alıştırma.
\end{proof}

\begin{prob}
\olref[fol][syn][unq]{prop:unique-atomic}'i ispatlayın (İpucu: Terimler
için \olref[fol][syn][unq]{lem:no-prefix}'in bir sürümünü ifade edip
ispatlayın.)
\end{prob}

\begin{prop}[Tek Türlü Okunabilirlik]
Her !!{formula} aşağıdaki koşullardan yalnızca birini sağlar.
\begin{enumerate}
\item $!A$ atomiktir.

\tagitem{prvNot}{$!A$, $\lnot !B$ biçimindedir.}{}

\tagitem{prvAnd}{$!A$, $(!B \land !C)$ biçimindedir.}{}

\tagitem{prvOr}{$!A$, $(!B \lor !C)$ biçimindedir.}{}

\tagitem{prvIf}{$!A$, $(!B \lif !C)$ biçimindedir.}{}

\tagitem{prvIff}{$!A$, $(!B \liff !C)$ biçimindedir.}{}

\tagitem{prvAll}{$!A$, $\lforall[x][!B]$ biçimindedir.}{}

\tagitem{prvEx}{$!A$, $\lexists[x][!B]$ biçimindedir.}{}
\end{enumerate}
Ayrıca her durumda $!B$ ya da $!B$ ile $!C$ tek türlü belirlenmiştir.
Bu, örneğin $!A$'nın hem
\iftag{prvIf}{$(!B \lif !C)$ hem de $(!B' \lif !C')$ biçiminde} {
    \iftag{prvOr}{$(!B \lor !C)$ hem de $(!B' \lor !C')$ biçiminde} {
      \iftag{prvAnd}{$(!B \land !C)$ hem de $(!B' \land !C')$ biçiminde}{}}}
olmasını sağlayan farklı $!B$, $!C$ ve $!B'$, $!C'$ çiftlerinin
bulunmadığı anlamına gelir.
\end{prop}

\begin{proof}
Oluşum kuralları, bir !!a{formula} atomik değilse bir açma
parantezi~(, \iftag{prvNot}{$\lnot$,}{} veya bir niceleyiciyle başlamasını
gerektirir. Öte yandan aşağıdaki simgelerden biriyle başlayan her
!!{formula} atomik olmalıdır: !!a{predicate}, bir değişken,
!!a{function}, !!a{constant}\iftag{prvFalse}{, $\lfalse$}{}\iftag{prvTrue}{, $\ltrue$}{}.

Dolayısıyla gerçekte yalnızca $!A$ hem $(!B \ast !C)$ hem de $(!B'
\mathbin{\ast'} !C')$ biçimindeyse $!B \ident !B'$, $!C \ident !C'$ ve
$\ast = {\ast'}$ olduğunu göstermeliyiz.

Hem $!A \ident (!B \ast !C)$ hem de $!A \ident (!B'
\mathbin{\ast'} !C')$ olduğunu varsayalım. Ya $!B \ident !B'$'dir ya
da değildir. İlki geçerliyse, aynı yerde başlayıp aynı uzunlukta olan
$!A$ alt dizileri oldukları için açıkça $\ast = {\ast'}$ ve
$!C \ident !C'$'dir. Öteki durumda $!B \not\ident !B'$'dir. $!B$ ve
$!B'$ aynı yerde başlayan $!A$ alt dizileri olduğundan, biri ötekinin
öz öneki olmalıdır. Fakat bu, \olref{lem:no-prefix} uyarınca olanaksızdır.
\end{proof}


